% Part: incompleteness
% Chapter: theories-computability
% Section: q-is-ce

\documentclass[../../../include/open-logic-section]{subfiles}

\begin{document}

\olfileid[pt]{inc}{tcp}{qce}

\olsection{$\Th{Q}$ é \printtoken{S}{c.e.}-completa}


\begin{thm}
$\Th{Q}$ é !!{c.e.}, mas não decidível. Na verdade, é um conjunto
  !!{c.e.} completo.
\end{thm}

\begin{proof}
Não é difícil ver que $\Th{Q}$ é !!{c.e.}, pois ela é o conjunto dos
(códigos de) sentenças $y$ tais que existe uma prova $x$ de $y$ em
$\Th{Q}$:
\[
Q = \Setabs{y}{\lexists[x][\Prf[\Th{Q}](x,y)]}.
\]
Mas sabemos que $\Prf[\Th{Q}](x,y)$ é computável (na verdade, recursiva
primitiva), e todo conjunto que pode ser escrito na forma acima é c.e.

Dizer que ela é um conjunto c.e.\ completo equivale a dizer que
$K \leq_m Q$, em que $K = \Setabs{x}{!A_x(x) \downarrow}$. Mostremos,
portanto, que $K$ é redutível a $\Th{Q}$. Como o predicado de Kleene
$T(e,x,s)$ é recursivo primitivo, ele é representável em $\Th{Q}$,
digamos, por $!A_T$. Então, para todo $x$, temos
\begin{align*}
x \in K & \lif \lexists[s][T(x,x,s)] \\
& \lif \lexists[s][(\Th{Q} \Proves !A_T(\num x,\num x, \num s))]
  \\
& \lif \Th{Q} \Proves \lexists[s][!A_T(\num x, \num x, s)].
\end{align*}
Reciprocamente, se $\Th{Q} \Proves
\lexists[s][!A_T(\num x, \num x, s)]$, então, de fato, para algum
número natural $n$, a fórmula $!A_T(\num x, \num x, \num n)$ deve ser
verdadeira. Ora, se $T(x,x,n)$ fosse falsa, $\Th{Q}$ provaria
$\lnot !A_T(\num x, \num x, \num n)$, pois $!A_T$ representa $T$.
Mas, nesse caso, $\Th{Q}$ provaria uma fórmula falsa, o que é uma
contradição. Logo, $T(x,x,n)$ deve ser verdadeira, o que implica
$!A_x(x) \downarrow$.

Em suma, para todo $x$, temos que $x$ pertence a $K$ se, e somente se,
$\Th{Q}$ prova $\lexists[s][T(\num x,\num x,s)]$. Assim, a função $f$
que leva $x$ a (um código para) a sentença $\lexists[s][T(\num x,
  \num x, s)]$ é uma redução de $K$ a $\Th{Q}$.
\end{proof}

\end{document}
